\chapter{쪽지 돌리기 놀이}

5장의 네 친구를 다시 불러 볼게요.
지우, 하준, 서연, 도윤이에요.
허락 지도는 그림~\ref{fig:four}와 같아요.

\section{놀이 규칙}

쪽지 한 장과 카드 100장이 있어요.
카드 85장에는 “화살표를 따라가”라고 적혀 있고,
카드 15장에는 “선생님께”라고 적혀 있어요.

\begin{enumerate}
\item 선생님이 쪽지를 아무 친구에게나 줘요.
\item 쪽지를 받은 친구는 카드를 한 장 뽑아요.
\item “화살표를 따라가”가 나오면, 자기 화살표가 가리키는 친구에게 쪽지를 넘겨요.
      화살표가 여러 개면 그중 하나를 골고루 골라요.
\item “선생님께”가 나오면 쪽지를 선생님께 드리고, 선생님은 다시 아무 친구에게나 줘요.
      화살표가 하나도 없는 친구도 쪽지를 선생님께 드려요.
\item 뽑은 카드는 다시 섞어 넣고, 2번부터 계속해요.
\end{enumerate}

이 놀이를 아주 오랫동안 해요.
천 번, 만 번, 백만 번이요.
그리고 쪽지가 누구 손에 몇 번 있었는지 세어요.

어떤 친구가 쪽지를 오래 갖게 될까요?
화살표를 많이 받는 친구는 쪽지가 자주 찾아와요.
그런데 그게 다가 아니에요.
쪽지를 자주 갖는 친구가 가리키는 친구도 쪽지를 자주 받아요.
쪽지가 많이 머무는 곳에서 흘러나온 쪽지가 그 친구에게 가기 때문이에요.

바로 이거예요.
5장 끝에서 우리는 “믿을 만한 친구가 보낸 화살표를 더 무겁게 치고 싶은데,
누가 믿을 만한지 모른다”는 꼬인 문제에 빠졌어요.
쪽지 놀이는 그 문제를 저절로 풀어 줘요.
미리 아무것도 알 필요가 없어요.
그냥 쪽지를 돌리면, 쪽지가 오래 머무는 친구가 곧 “믿을 만한 친구에게서 많이 믿음을 받은 친구”가 돼요.
이렇게 정한 점수를 \textbf{페이지랭크}(PageRank)라고 불러요.

\section{“선생님께” 카드는 왜 있을까요?}

카드가 모두 “화살표를 따라가”라면 어떻게 될까요?
화살표가 없는 친구에게 가면 쪽지가 멈춰 버려요.
또 두 친구가 서로만 가리키고 있으면 쪽지가 그 둘 사이에서만 영원히 왔다 갔다 해요.
“선생님께” 카드는 가끔 쪽지를 교실 아무 데나 다시 던져 넣어서 이런 일을 막아 줘요.
85와 15라는 숫자는 페이지랭크를 처음 만든 사람들이 쓴 값이고,
제 논문에서도 같은 값을 썼어요.

\section{백만 번 대신 물 붓기}

놀이를 정말 백만 번 하려면 너무 오래 걸려요.
다행히 같은 답을 훨씬 빨리 얻는 방법이 있어요.
쪽지 대신 \emph{물}을 쓰는 거예요.

\begin{enumerate}
\item 처음에 네 친구가 물을 25컵씩, 모두 100컵 갖고 시작해요.
\item 한 차례마다 모든 친구가 자기 물을 한꺼번에 내보내요.
      85\%는 자기 화살표들로 똑같이 나누어 흘려보내고,
      15\%는 선생님 양동이에 부어요.
\item 선생님은 양동이의 물을 네 친구에게 똑같이 나누어 줘요.
\item 이 차례를 계속 되풀이해요.
\end{enumerate}

첫 차례를 손으로 계산해 볼게요.
선생님 양동이에는 모두의 물 100컵 가운데 15\%인 15컵이 모이고, 한 명당 3.75컵씩 돌아가요.
지우는 25컵의 85\%인 21.25컵을 서연이에게 보내요.
도윤이도 21.25컵을 서연이에게 보내요.
하준이는 화살표가 두 개라서 21.25컵을 반으로 나누어 서연이와 도윤이에게 10.625컵씩 보내요.
서연이는 21.25컵을 하준이에게 보내요.
그러면 첫 차례가 끝났을 때

\begin{itemize}
\item 지우: 선생님께 받은 3.75컵뿐 → 3.8컵
\item 하준: 3.75 + 21.25 → 25.0컵
\item 서연: 3.75 + 21.25 + 21.25 + 10.625 → 56.9컵
\item 도윤: 3.75 + 10.625 → 14.4컵
\end{itemize}

이 돼요.
이것을 계속하면 이렇게 바뀌어 가요.

\begin{table}[h]
\centering
\small
\begin{tabular}{@{}lrrrr@{}}
\toprule
차례 & 지우 & 하준 & 서연 & 도윤 \\
\midrule
처음 & 25.0 & 25.0 & 25.0 & 25.0 \\
1 & 3.8 & 25.0 & 56.9 & 14.4 \\
2 & 3.8 & 52.1 & 29.8 & 14.4 \\
3 & 3.8 & 29.1 & 41.3 & 25.9 \\
4 & 3.8 & 38.9 & 41.3 & 16.1 \\
5 & 3.8 & 38.9 & 37.1 & 20.3 \\
\multicolumn{1}{c}{$\vdots$} & & & & \\
13 이후 & \textbf{3.8} & \textbf{37.3} & \textbf{39.4} & \textbf{19.6} \\
\bottomrule
\end{tabular}
\caption{물 붓기를 되풀이하면 물의 양이 출렁이다가 한곳에 자리를 잡아요(단위: 컵, 모두 더하면 100).}
\end{table}

처음 몇 차례는 물이 이리저리 출렁여요.
그러다 열세 번째 차례쯤부터는 거의 움직이지 않아요.
이때의 물 양이 쪽지 놀이를 백만 번 했을 때 쪽지가 머무는 비율과 같아요.
이것이 페이지랭크 점수예요.

\section{결과를 읽어 봐요}

\begin{itemize}
\item \textbf{서연}이가 39.4로 1등이에요. 화살표를 3개나 받았으니 그럴 만하지요.
\item \textbf{하준}이가 37.3으로 서연이와 거의 같아요. 받은 화살표는 1개뿐인데요!
      그 1개가 1등인 서연이에게서 왔고, 서연이는 물을 \emph{모두} 하준이에게 보내기 때문이에요.
\item \textbf{도윤}이는 19.6이에요. 받은 화살표 1개가 하준이에게서 오지만, 하준이는 물을 둘로 나누어 보내요.
\item \textbf{지우}는 3.8, 선생님께 받는 몫뿐이에요. 아무도 지우를 가리키지 않으니까요.
\end{itemize}

받은 화살표를 세었을 때는 하준이와 도윤이가 1개씩으로 같았어요.
페이지랭크는 둘을 구별해요.
“누가 보낸 화살표인가”를 따지기 때문이에요.

\section{가짜 친구 작전, 다시}

이제 5장의 도윤이 작전을 페이지랭크로 다시 계산해 볼게요.
가짜 통 세 개가 생겨서 통이 모두 일곱 개가 되었어요.
물 100컵을 일곱 통이 나누어 시작해요.

가짜 통에는 아무도 화살표를 보내지 않아요.
그래서 가짜 통이 가진 물은 선생님이 나누어 주는 몫, 한 통에 2.1컵뿐이에요.
도윤이에게 흘러가는 물도 거기서 조금 나오는 것이 전부예요.

\begin{center}
\small
\begin{tabular}{@{}lccc@{}}
\toprule
& 받은 화살표 수 & 페이지랭크 & 페이지랭크 순위 \\
\midrule
서연 & 3 & 36.5 & 1등 \\
하준 & 1 & 33.2 & 2등 \\
도윤 & \textbf{4} & 21.7 & 3등 \\
지우 & 0 & 2.1 & \\
가짜 통 세 개 & 0 & 2.1씩 & \\
\bottomrule
\end{tabular}
\end{center}

받은 화살표 수로는 도윤이가 1등이 되었지만, 페이지랭크로는 여전히 3등이에요.
아무도 믿지 않는 통이 보낸 화살표는 무게가 아주 가볍기 때문이에요.
다만 완벽하지는 않아요.
도윤이의 물은 19.6에서 21.7로 조금 늘었어요.
가짜 통을 아주 많이 만들면 효과가 쌓일 수 있어요.
페이지랭크는 속임수를 \emph{어렵게} 만들 뿐, 완전히 막지는 못해요.

\begin{deeper}[페이지랭크의 역사와 식]
페이지랭크는 1998년 미국 스탠퍼드 대학교 대학원생이던 래리 페이지(Larry Page)와 세르게이 브린(Sergey Brin)이 웹 페이지의 순위를 매기려고 만들었어요.
웹 페이지가 점이고, 한 페이지가 다른 페이지로 거는 링크가 화살표예요.
이름의 “페이지”는 웹 페이지와 래리 페이지의 성을 함께 뜻해요.
두 사람은 이 생각으로 구글을 세웠어요.

쪽지가 화살표를 따라 이리저리 옮겨 다니는 놀이는 러시아의 수학자 안드레이 마르코프가 1906년에 연구하기 시작한 \textbf{마르코프 연쇄}의 한 예예요.

식으로 쓰면, 점 $v$의 점수 $\mathrm{PR}(v)$는
\[
\mathrm{PR}(v) = \frac{1-d}{N} + d \sum_{u \to v} \frac{\mathrm{PR}(u)}{\text{($u$의 화살표 수)}}
\]
이에요. $d=0.85$는 “화살표를 따라가” 카드의 비율, $N$은 점의 수예요.
화살표에 무게가 있으면 나누는 몫을 무게에 비례하게 바꿔요(7장).
물 붓기처럼 되풀이해서 푸는 방법을 \textbf{거듭제곱법}(power iteration)이라고 해요.
논문에서는 변화가 $10^{-8}$보다 작아지거나 300번을 채우면 멈추었어요.
실제 지갑 5,521개로 계산했을 때 허락 지도는 37번, 송금 지도는 100번 만에 멈추었어요.
\end{deeper}

\begin{learned}
\begin{itemize}
\item 쪽지 돌리기 놀이에서 쪽지가 오래 머무는 친구가 페이지랭크 점수가 높아요.
\item 페이지랭크는 “믿을 만한 친구가 보낸 화살표를 더 무겁게”를 저절로 해 줘요.
\item 물 붓기를 되풀이하면 놀이를 백만 번 하지 않고도 점수를 구할 수 있어요.
\item 가짜 통의 화살표는 가볍지만, 속임수를 완전히 막지는 못해요.
\end{itemize}
\end{learned}
